\chapter{Лабораторная работа №6. Создание WLAN.}
\label{ch:chapter 6}

\section{Цели}

Лабораторная работа помогает получить практические навыки по изучению
следующих тем:

\begin{itemize}

\item Процедура аутентификации точек доступа

\item Процедура настройки профилей WLAN

\item Процесс настройки основных параметров WLAN

\end{itemize}

\section{Топология сети}

\includegraphics[width=0.5\textwidth]{lb6_1.png} 
\captionof{figure}{Топология сети для лабораторной работы №6}

\section{План работы}

\begin{itemize}

\item 1. Настройка подключения к проводной сети.

\item 2. Настройка точек доступа и перевод их в режим онлайн.

(1) Создание групп точек доступа и добавление точек доступа с одинаковой конфигурацией в одну группу для унифицированной настройки

(2) Настройка системных параметров контроллера доступа, включая код страны и интерфейс-источник, используемый контроллером для связи с точками доступа.

(3) Настройка режима аутентификации AP и импорт AP для выхода точек доступа в сеть.

\item 3. Настройка параметров сервисов WLAN и передача конфигурации точкам доступа, чтобы обеспечить доступ STA к WLAN.

\end{itemize}

\section{Процедура конфигурирования}

К сожалению, ни один из AP, представленных в виртуальном окружении не работает, а потому невозможно выполнить настройку WLAN, поэтому далее будет представлена настройка Switch.

\textbf{Шаг 1.} Настроим основные параметры устройства.

Отключим ненужные порты между S1 и AC:

\includegraphics[width=0.5\textwidth]{lb6_2.png} 

Включим функцию PoE (Power over Ethernet) на портах S3 и S4, подключенных к AP:

\includegraphics[width=0.5\textwidth]{lb6_3.png} 

\textbf{Шаг 2.} Настроим параметры проводной сети.

Настроим VLAN:

\includegraphics[width=0.5\textwidth]{lb6_4.png} 

\includegraphics[width=0.5\textwidth]{lb6_5.png} 

\includegraphics[width=0.5\textwidth]{lb6_6.png} 

\includegraphics[width=0.5\textwidth]{lb6_7.png} 

Настроим IP-адреса интерфейсов:

\includegraphics[width=0.5\textwidth]{lb6_8.png} 

Настроим DHCP:

\includegraphics[width=0.5\textwidth]{lb6_9.png} 

\textbf{Шаг 3.} Настроим параметры точек доступа для выхода в сеть.

В силу вышеописанной проблемы с AP, я просто опишу, как именно настроить параметры точек доступа.

Создадим группу AP и назовем ее ap-group1.

\texttt{[AC]wlan}

\texttt{[AC-wlan-view]ap-group name ap-group1}

\texttt{[AC-wlan-ap-group-ap-group1]quit}

Создадим профиль регулирующего домена и настроим код страны AC в профиле.

\texttt{[AC]wlan}

\texttt{[AC-wlan-view]regulatory-domain-profile name default}

\texttt{[AC-wlan-regulate-domain-default]country-code cn}

\texttt{[AC-wlan-regulate-domain-default]quit}

Установим привязку профиля регулирующего домена к группе AP.

\texttt{[AC]wlan}

\texttt{[AC-wlan-view]ap-group name ap-group1}

\texttt{[AC-wlan-ap-group-ap-group1]regulatory-domain-profile default}

\texttt{[AC-wlan-ap-group-ap-group1]quit}

Укажем интерфейс-источник на AC для установления туннелей CAPWAP.

\texttt{[AC]capwap source interface Vlanif 100}

Импортируем точки доступа в AC и добавим их в группу AP с именем ap-group1.

\texttt{[AC]wlan}

\texttt{[AC-wlan-view]ap auth-mode mac-auth}

\texttt{[AC-wlan-view]ap-id 0 ap-mac 60F1-8A9C-2B40}

\texttt{[AC-wlan-ap-0]ap-name ap1}

\texttt{[AC-wlan-ap-0]ap-group ap-group1}

\texttt{[AC-wlan-ap-0]quit}

\texttt{[AC-wlan-view]ap-id 1 ap-mac B4FB-F8B7-DE40}

\texttt{[AC-wlan-ap-1]ap-name ap2}

\texttt{[AC-wlan-ap-1]ap-group ap-group1}

\texttt{[AC-wlan-ap-1]quit}

\textbf{Шаг 4.} Настроим параметры сервисов WLAN.

Создадим профиль безопасности HCIA-WLAN и настроим политику безопасности.

\texttt{[AC-wlan-view]security-profile name HCIA-WLAN}

\texttt{[AC-wlan-sec-prof-HCIA-WLAN]security wpa-wpa2 psk pass-phrase HCIA-Datacom aes}

\texttt{[AC-wlan-sec-prof-HCIA-WLAN]quit}

Создадим профиль SSID HCIA-WLAN и зададим имя SSID HCIA-WLAN.

\texttt{[AC]wlan}

\texttt{[AC-wlan-view]ssid-profile name HCIA-WLAN}

\texttt{[AC-wlan-ssid-prof-HCIA-WLAN]ssid HCIA-WLAN}

\texttt{[AC-wlan-ssid-prof-HCIA-WLAN]quit}

Создадим профиль VAP HCIA-WLAN, настроим режим передачи данных и сервисную VLAN и применим профиль безопасности и профиль SSID к профилю VAP.

\texttt{[AC]wlan}

\texttt{[AC-wlan-view]vap-profile name HCIA-WLAN}

\texttt{[AC-wlan-vap-prof-HCIA-WLAN]forward-mode direct-forward}

\texttt{[AC-wlan-vap-prof-HCIA-WLAN]service-vlan vlan-id 101}

\texttt{[AC-wlan-vap-prof-HCIA-WLAN]security-profile HCIA-WLAN}

\texttt{[AC-wlan-vap-prof-HCIA-WLAN]ssid-profile HCIA-WLAN}

\texttt{[AC-wlan-vap-prof-HCIA-WLAN]quit}

Установим привязку профиля VAP к группе AP и применим конфигурацию профиля VAP HCIA-WLAN к радиомодулю 0 и радиомодулю 1 точек доступа в группе AP.

\texttt{[AC]wlan}

\texttt{[AC-wlan-view]ap-group name ap-group1}

\texttt{[AC-wlan-ap-group-ap-group1]vap-profile HCIA-WLAN wlan 1 radio all}

\texttt{[AC-wlan-ap-group-ap-group1]quit}

\section*{Вопросы по лабораторной работе WLAN}

\subsection*{Вопрос 1}
\textbf{Какое влияние на доступ STA к S1 в текущей сети будут оказывать настройки порта GigabitEthernet0/0/10 контроллера доступа, запрещающие прохождение пакетов из VLAN 101? Почему? Что изменится, если будет использоваться туннельная передача?}

\textbf{Ответ:}

При использовании \textbf{прямой передачи (direct-forward)} настройки порта GigabitEthernet0/0/10 контроллера доступа, запрещающие прохождение пакетов из VLAN 101, \textbf{не будут оказывать никакого влияния} на доступ STA к S1.

\textbf{Почему:} 
При прямой передаче данных пользовательский трафик от STA (станций) не проходит через контроллер доступа (AC). Трафик передается от точки доступа (AP) напрямую на коммутатор S3/S4, затем на S1, минуя AC. AC участвует только в управлении AP и не обрабатывает пользовательские данные. Следовательно, запрет на порту AC не влияет на прохождение пользовательского трафика.

\textbf{Что изменится при туннельной передаче:}
При использовании туннельной передачи (туннелирование CAPWAP) пользовательский трафик инкапсулируется в туннель CAPWAP и передается от AP к AC, а затем AC декапсулирует его и отправляет в проводную сеть. В этом случае, если порт AC запрещает прохождение пакетов из VLAN 101, то STA \textbf{не смогут получить доступ к S1}, так как трафик будет блокироваться на порту контроллера доступа.

\subsection*{Вопрос 2}
\textbf{Какие операции необходимо выполнить на AC, чтобы назначить STA, подключенные к AP1 и AP2, к разным VLAN?}

\textbf{Ответ:}

Для назначения STA, подключенных к разным точкам доступа (AP1 и AP2), в разные VLAN, необходимо выполнить следующие операции на контроллере доступа (AC):

\begin{enumerate}
\item \textbf{Создать разные профили VAP (Virtual Access Point) для разных AP.}
    
    \texttt{[AC-wlan-view]vap-profile name WLAN-AP1}

    \texttt{[AC-wlan-vap-prof-WLAN-AP1]service-vlan vlan-id 101}
    
    \texttt{[AC-wlan-vap-prof-WLAN-AP1]quit}
    
    \texttt{[AC-wlan-view]vap-profile name WLAN-AP2}

    \texttt{[AC-wlan-vap-prof-WLAN-AP2]service-vlan vlan-id 102}

    \texttt{[AC-wlan-vap-prof-WLAN-AP2]quit}

\item \textbf{Создать разные группы AP (или использовать существующие группы).}
    
    \texttt{[AC-wlan-view]ap-group name group-AP1}

    \texttt{[AC-wlan-view]ap-group name group-AP2}

\item \textbf{Привязать соответствующие профили VAP к радиомодулям каждой группы AP.}
    
    \texttt{[AC-wlan-view]ap-group name group-AP1}

    \texttt{[AC-wlan-ap-group-group-AP1]vap-profile WLAN-AP1 wlan 1 radio all}

    \texttt{[AC-wlan-ap-group-group-AP1]quit}
    
    \texttt{[AC-wlan-view]ap-group name group-AP2}

    \texttt{[AC-wlan-ap-group-group-AP2]vap-profile WLAN-AP2 wlan 1 radio all}

    \texttt{[AC-wlan-ap-group-group-AP2]quit}

\item \textbf{Назначить каждую точку доступа в соответствующую группу AP.}
    
    \texttt{[AC-wlan-view]ap-id 0}

    \texttt{[AC-wlan-ap-0]ap-group group-AP1}

    \texttt{[AC-wlan-ap-0]quit}
    
    \texttt{[AC-wlan-view]ap-id 1}

    \texttt{[AC-wlan-ap-1]ap-group group-AP2}

    \texttt{[AC-wlan-ap-1]quit}
\end{enumerate}
